%=============================================================================
\section{Digital Signal Processing}
%=============================================================================

You might be wondering: ``I signed up for control theory, why are we starting with signal processing?'' Great question. Here is the answer: \textbf{garbage in, garbage out.}

Your controller is only as good as the data it receives. If your IMU reading is bouncing around like a caffeinated squirrel, no amount of PID tuning will save you. So before we learn to control systems, we learn to clean up their signals.

\subsection{Understanding the Frequency Domain}

In everyday life, we think about signals in the \textbf{time domain}---voltage changes over time, position changes over time. This is natural and intuitive. But it turns out that many problems become \textit{much} easier if we look at signals from a different angle: the \textbf{frequency domain}.

The core idea is beautifully simple: \textit{any signal can be decomposed into a sum of sinusoids of different frequencies.} Your noisy gyroscope reading? It is actually a clean low-frequency rotation signal plus a bunch of high-frequency noise signals all added together. If we can identify which frequencies are ``signal'' and which are ``noise,'' we can surgically remove the noise.

Think of it like an audio equalizer on a music player. The bass, mids, and treble are all mixed together in the time-domain waveform. But the equalizer lets you boost or cut specific frequency ranges independently. That is exactly what we do with filters in control systems.

\textbf{Key intuition:} Slow, meaningful changes in your system (like a chassis turning) happen at low frequencies. Fast, meaningless changes (like electrical noise or vibration) happen at high frequencies. The frequency domain lets us tell them apart.

\subsection{Laplace Transform and Fourier Transform}

These two transforms are your tickets from the time domain to the frequency domain. They are closely related, but they serve slightly different purposes.

\subsubsection{Laplace Transform}

The Laplace transform takes a time-domain function $f(t)$ and converts it into a function of the complex variable $s$:

\begin{equation}
F(s) = \mathcal{L}\{f(t)\} = \int_0^{\infty} f(t) e^{-st} \, dt
\end{equation}

where $s = \sigma + j\omega$ is a complex number. The real part $\sigma$ captures exponential growth or decay, and the imaginary part $j\omega$ captures oscillation.

\textbf{Why do we care?} Because the Laplace transform turns differential equations into algebraic equations. Instead of solving
\[
m\ddot{x} + b\dot{x} + kx = F(t),
\]
(which requires calculus), we can solve
\[
(ms^2 + bs + k)X(s) = F(s),
\]
(which is just algebra). Derivatives become multiplication by $s$. Integration becomes division by $s$. Life becomes easier.

The most important Laplace transforms you will use daily:

\begin{center}
\begin{tabular}{lll}
\toprule
\textbf{Time domain} & \textbf{$s$-domain} & \textbf{When you see it} \\
\midrule
$1$ (step input) & $\frac{1}{s}$ & Commanding a new setpoint \\
$e^{-at}$ & $\frac{1}{s+a}$ & First-order system response \\
$\sin(\omega t)$ & $\frac{\omega}{s^2 + \omega^2}$ & Oscillation, vibration \\
$\dot{f}(t)$ & $sF(s) - f(0)$ & Derivative in the controller \\
$\int f(t)\,dt$ & $\frac{F(s)}{s}$ & Integral in the controller \\
\bottomrule
\end{tabular}
\end{center}

\textbf{Pro tip:} You do not need to memorize the integral formula. You need to understand that $s$ means ``derivative'' and $1/s$ means ``integral.'' That single insight will carry you through 80\% of control theory.

\textbf{What is $e^{-st}$ actually doing?} The kernel $e^{-st} = e^{-\sigma t} e^{-j\omega t}$ is a ``probe signal'' that asks two questions simultaneously: (1) $e^{-\sigma t}$: ``Does the signal grow or decay at rate $\sigma$?'' (2) $e^{-j\omega t}$: ``How much of frequency $\omega$ is present?'' By integrating $f(t) \cdot e^{-st}$ over all time, we compute the ``correlation'' between the signal and every possible exponential-sinusoid. The result $F(s)$ is a map of the entire complex plane, telling us exactly how much of each exponential-oscillatory mode exists in the signal.

\textbf{Common Laplace transform pairs --- worked out:}

The \textbf{unit step} $u(t) = 1$ for $t \geq 0$:
\[
\mathcal{L}\{1\} = \int_0^{\infty} e^{-st} \, dt = \left[-\frac{e^{-st}}{s}\right]_0^{\infty} = \frac{1}{s}
\]

The \textbf{exponential decay} $e^{-at}$:
\[
\mathcal{L}\{e^{-at}\} = \int_0^{\infty} e^{-at} e^{-st} \, dt = \int_0^{\infty} e^{-(s+a)t} \, dt = \frac{1}{s+a}
\]
This is why a first-order system with pole at $s = -a$ has a time response of $e^{-at}$---the pole location \textit{is} the decay rate.

The \textbf{ramp} $t \cdot u(t)$:
\[
\mathcal{L}\{t\} = \frac{1}{s^2}
\]
Notice: multiplying by $t$ in time corresponds to differentiating with respect to $s$ and flipping the sign: $\mathcal{L}\{t \cdot f(t)\} = -\frac{d}{ds}F(s)$. Since $\mathcal{L}\{1\} = 1/s$, we get $\mathcal{L}\{t\} = -\frac{d}{ds}\frac{1}{s} = \frac{1}{s^2}$.

\subsubsection{Fourier Transform}

The Fourier transform is the Laplace transform's well-behaved cousin. It is what you get when you set $s = j\omega$ (i.e., $\sigma = 0$):

\begin{equation}
F(j\omega) = \int_{-\infty}^{\infty} f(t) e^{-j\omega t} \, dt
\end{equation}

While the Laplace transform works with the full complex plane (and can handle unstable, growing signals), the Fourier transform lives strictly on the imaginary axis. It tells you: ``How much of each frequency is present in this signal?''

\textbf{When to use which:}
\begin{itemize}
    \item \textbf{Laplace transform}: For analyzing and designing control systems (transfer functions, stability).
    \item \textbf{Fourier transform}: For analyzing signals (noise, vibration, frequency content).
\end{itemize}

In practice, for digital systems we often use the \textbf{Discrete Fourier Transform (DFT)}, computed efficiently via the \textbf{Fast Fourier Transform (FFT)} algorithm. If you ever want to see what frequencies are hiding in your sensor data, just throw it into an FFT. MATLAB's \texttt{fft()} or Python's \texttt{numpy.fft.fft()} will do the job in one line.

\subsubsection{Frequency Spectrum}

When you apply the Fourier transform to a signal, you get its \textbf{frequency spectrum}---a plot showing the amplitude (and phase) of each frequency component.

Reading a frequency spectrum is like reading an X-ray of your signal:
\begin{itemize}
    \item A tall peak at 50 Hz? That is probably power line interference.
    \item A broad hump below 10 Hz? That is likely your actual signal of interest.
    \item Random fuzz spread across all frequencies? That is white noise.
\end{itemize}

\textbf{Practical example:} Suppose your gimbal encoder gives a noisy velocity reading. You record 1 second of data, compute the FFT, and see a clean peak at 3 Hz (your gimbal is tracking a target moving at 3 Hz) plus a mess of noise above 50 Hz. Now you know: design a low-pass filter with a cutoff around 20--30 Hz, and your signal will be clean.

The frequency spectrum has two parts:
\begin{itemize}
    \item \textbf{Magnitude spectrum} $|F(j\omega)|$: How strong each frequency component is.
    \item \textbf{Phase spectrum} $\angle F(j\omega)$: The time shift of each frequency component.
\end{itemize}

For filter design, we mostly care about the magnitude spectrum. Phase becomes important later when we analyze control loop stability (phase margin).

\begin{figure}[H]
\centering
\includegraphics[width=\textwidth]{frequency_spectrum.pdf}
\caption{Top: a noisy time-domain signal looks hopeless. Bottom: the FFT instantly reveals a 5 Hz signal and a 42 Hz resonance hiding in the noise.}
\end{figure}

\vspace{0.5em}
\begin{mdframed}[linewidth=1pt, innertopmargin=10pt, innerbottommargin=10pt]
\textbf{Case Study: Debugging a Vibrating Chassis with FFT}

\vspace{0.3em}
True story pattern from competition robotics: Your mecanum chassis drives fine at low speed but develops a nasty vibration at high speed. The whole frame buzzes. Everyone suspects a loose screw. You tighten everything. Still vibrates.

\textbf{The diagnosis:} Log the IMU accelerometer data for 5 seconds while driving at the problematic speed. Run an FFT in Python:

\begin{verbatim}
    import numpy as np
    freqs = np.fft.rfftfreq(len(data), d=1.0/sample_rate)
    spectrum = np.abs(np.fft.rfft(data))
    plt.plot(freqs, spectrum)
\end{verbatim}

You see a massive spike at exactly 42 Hz. Interesting. You calculate: at high speed, the mecanum wheel rollers contact the ground at\ldots about 42 Hz. The roller contact frequency matches a structural resonance of the chassis frame.

\textbf{The fix options:}
\begin{itemize}[nosep]
    \item Add a \textbf{notch filter} (band-reject) at 42 Hz in the control loop to prevent the controller from exciting the resonance.
    \item Add physical damping (rubber mounts) to shift the resonance frequency.
    \item Stiffen the frame to move the resonance above the operating range.
\end{itemize}

\textbf{The lesson:} Without the FFT, you would have spent hours tightening screws. With the FFT, the root cause was obvious in 10 seconds. \textbf{The frequency domain is a debugging superpower.} Whenever something oscillates and you do not know why, record the data and look at the spectrum.
\end{mdframed}

\subsection{Frequency-Based Filter Design}

Now that we can see signals in the frequency domain, we can design filters to keep the frequencies we want and reject the ones we do not.

A filter is simply a system that has different gains at different frequencies. We describe this using the filter's \textbf{frequency response}---a plot of gain vs.\ frequency (called a \textbf{Bode magnitude plot}).

\subsubsection{Low-Pass Filter}

The most common filter in robotics. It passes low frequencies and attenuates high frequencies.

\textbf{First-order low-pass filter:}
\begin{equation}
H(s) = \frac{\omega_c}{s + \omega_c}
\end{equation}

where $\omega_c = 2\pi f_c$ is the cutoff frequency in rad/s. At frequencies below $\omega_c$, the gain is approximately 1 (signal passes through). Above $\omega_c$, the gain rolls off at $-20$ dB/decade (signal gets attenuated).

\textbf{In the time domain}, this is equivalent to an exponential moving average:
\begin{equation}
y[n] = \alpha \cdot x[n] + (1 - \alpha) \cdot y[n-1]
\end{equation}

where $\alpha = \frac{T_s}{T_s + \frac{1}{2\pi f_c}}$ and $T_s$ is the sampling period.

\textbf{When to use:} Smoothing noisy sensor readings (gyroscope, accelerometer, motor current). This is the filter you will use 90\% of the time.

\textbf{Gotcha:} A low-pass filter introduces \textbf{phase lag}---the output is delayed relative to the input. The more aggressively you filter (lower cutoff), the more lag you introduce. In a fast control loop, this lag can destabilize your system. There is always a trade-off between smooth signals and fast response.

\begin{figure}[H]
\centering
\includegraphics[width=\textwidth]{lpf_response.pdf}
\caption{Left: LPF step response --- lower cutoff means slower rise but smoother output. Right: filtering a noisy 3 Hz signal --- the LPF recovers the true signal while removing high-frequency noise.}
\end{figure}

\vspace{0.5em}
\begin{mdframed}[linewidth=1pt, innertopmargin=10pt, innerbottommargin=10pt]
\textbf{Case Study: Why a Low-Pass Filter Gives You ``Soft Start'' for Free}

\vspace{0.3em}
Suppose you command your motor to jump from 0 to 5000 RPM instantly (a step input). Without any filter, the setpoint changes discontinuously---the motor driver sees a sudden full-voltage command, the gears slam, the chassis lurches, and your teammates spill their coffee.

Now insert a first-order LPF on the setpoint: $r_{\text{filtered}}(s) = \frac{\omega_c}{s + \omega_c} \cdot R(s)$. What happens?

The step input in the $s$-domain is $R(s) = \frac{5000}{s}$. After the filter:
\[
R_{\text{filtered}}(s) = \frac{5000 \, \omega_c}{s(s + \omega_c)}
\]

Taking the inverse Laplace transform:
\[
r_{\text{filtered}}(t) = 5000 \left(1 - e^{-\omega_c t}\right)
\]

The step has been transformed into a \textbf{smooth exponential ramp}. The time constant is $\tau = 1/\omega_c$. After $3\tau$, the output reaches 95\% of the target; after $5\tau$, it is at 99\%.

\textbf{Why this matters on your robot:}
\begin{itemize}[nosep]
    \item \textbf{Mechanical protection:} Gears and belts have torque limits. A step command can strip teeth or snap belts. The LPF limits the rate of change, keeping forces within safe bounds.
    \item \textbf{Current limiting:} A sudden voltage step causes a huge inrush current ($I = V/R$ when back-EMF is still zero). The LPF ramps up the voltage gradually, keeping current manageable.
    \item \textbf{Tunable aggressiveness:} Want a faster ramp? Increase $\omega_c$ (higher cutoff). Want a gentler ramp? Decrease $\omega_c$. One parameter controls the entire behavior.
\end{itemize}

\textbf{Comparison with a linear ramp:} You could also ramp the setpoint linearly (e.g., increase by 1000 RPM/s). But the LPF is better because it also smooths out \textit{any} discontinuity in the setpoint---not just the initial step, but also sudden changes during operation. It is a ``set it and forget it'' solution.

\textbf{The code is trivial:}
\begin{verbatim}
    setpoint_filtered += alpha * (setpoint_raw - setpoint_filtered);
\end{verbatim}

One line of code. No lookup tables, no state machines. That is the beauty of the LPF as a soft-start mechanism.
\end{mdframed}

\begin{figure}[H]
\centering
\includegraphics[width=\textwidth]{lpf_soft_start.pdf}
\caption{LPF soft start: a single parameter $\tau$ controls how aggressively the motor ramps up. The step command becomes a smooth exponential curve.}
\end{figure}

\subsubsection{Band-Pass Filter}

Passes a specific range of frequencies and attenuates everything else. Think of it as a low-pass and a high-pass filter combined.

\begin{equation}
H(s) = \frac{\frac{\omega_0}{Q} s}{s^2 + \frac{\omega_0}{Q} s + \omega_0^2}
\end{equation}

where $\omega_0$ is the center frequency and $Q$ is the quality factor (higher $Q$ = narrower band).

\textbf{When to use:} Extracting a specific frequency component from a signal. For example, if you know your target oscillates at a specific frequency and you want to track it while rejecting everything else.

In robotics, band-pass filters are less common than low-pass filters, but they are useful for vibration analysis and resonance detection.

\subsubsection{High-Pass Filter}

Passes high frequencies, attenuates low frequencies. It is the opposite of a low-pass filter.

\textbf{First-order high-pass filter:}
\begin{equation}
H(s) = \frac{s}{s + \omega_c}
\end{equation}

\textbf{When to use:} Removing DC offset or slow drift from a signal. For example, if your accelerometer has a slowly drifting bias, a high-pass filter can remove it while keeping the fast acceleration data intact.

\textbf{Practical note:} A complementary filter (commonly used for IMU sensor fusion) is just a low-pass filter on one sensor and a high-pass filter on another, with their cutoff frequencies matched so the gains add up to 1. The accelerometer gets low-passed (trust it for slow/static orientation) and the gyroscope gets high-passed (trust it for fast rotations). Simple, elegant, and surprisingly effective.

\subsubsection{Filter Type Selection and Parameter Tuning}

Choosing the right filter comes down to answering three questions:
\begin{enumerate}
    \item \textbf{What frequencies do I want to keep?} This determines the filter type (low-pass, high-pass, band-pass).
    \item \textbf{Where is the boundary?} This determines the cutoff frequency $f_c$.
    \item \textbf{How sharp should the transition be?} This determines the filter order.
\end{enumerate}

\textbf{Rules of thumb:}
\begin{itemize}
    \item Start with a first-order low-pass filter. It is simple, stable, and good enough for most applications.
    \item Set the cutoff frequency to 5--10$\times$ the bandwidth of your signal of interest. Too low and you lose signal; too high and you keep noise.
    \item If a first-order filter is not enough, try second-order (Butterworth for flat passband, Chebyshev for sharper rolloff).
    \item \textbf{Never} design a filter without looking at the frequency spectrum of your actual data first. Guessing cutoff frequencies is a recipe for frustration.
    \item When in doubt, plot the filtered output alongside the raw signal and check if the result makes physical sense.
\end{itemize}

\subsubsection{Digital Filter Structures: FIR vs.\ IIR}

So far we have described filters using continuous-time transfer functions $H(s)$.
Once you discretize them (covered in detail in Section~5), you end up with one of
two fundamental digital filter structures: \textbf{FIR} and \textbf{IIR}.
Understanding the difference is essential for choosing the right implementation on
your microcontroller.

\paragraph{FIR (Finite Impulse Response).}
An FIR filter computes the output as a weighted sum of the current and past input
samples only---no feedback:

\begin{equation}
y[n] = \sum_{k=0}^{M} b_k \, x[n-k]
\end{equation}

In the $z$-domain:
\begin{equation}
H(z) = \sum_{k=0}^{M} b_k \, z^{-k}
\end{equation}

FIR filters have \textbf{only zeros} (no poles except trivially at $z = 0$), which
gives them two important properties:
\begin{itemize}
    \item \textbf{Always stable}---no feedback means no risk of runaway oscillation.
    \item \textbf{Linear phase} is achievable by making the coefficients symmetric
          ($b_k = b_{M-k}$). This means every frequency component is delayed by the
          same amount, so the signal shape is perfectly preserved.
\end{itemize}

The price you pay is \textbf{filter order}: to get a sharp cutoff you may need
$M = 50$--$200$ taps, which means storing that many past samples and performing
that many multiplications per output sample.

\textbf{Common FIR design methods:}
\begin{itemize}
    \item \textbf{Window method}---truncate the ideal impulse response and apply a
          window function (Hamming, Blackman, Kaiser) to reduce spectral leakage.
    \item \textbf{Frequency sampling}---specify the desired magnitude response at
          DFT sample points and inverse-transform.
    \item \textbf{Parks--McClellan (Remez)}---an optimal algorithm that minimizes
          the maximum approximation error (equiripple design).
\end{itemize}

\paragraph{IIR (Infinite Impulse Response).}
An IIR filter uses \textit{both} past inputs and past outputs---it has a feedback path:

\begin{equation}
y[n] = \sum_{k=0}^{M} b_k \, x[n-k] \;-\; \sum_{k=1}^{N} a_k \, y[n-k]
\end{equation}

In the $z$-domain:
\begin{equation}
H(z) = \frac{\displaystyle\sum_{k=0}^{M} b_k \, z^{-k}}
            {1 + \displaystyle\sum_{k=1}^{N} a_k \, z^{-k}}
     = \frac{B(z)}{A(z)}
\end{equation}

IIR filters have \textbf{both poles and zeros}. The poles enable sharp frequency
selectivity with very low order (a 4th-order IIR can outperform a 50th-order FIR),
but they come with risks:
\begin{itemize}
    \item \textbf{Potentially unstable}---if any pole lies outside the unit circle
          ($|z| > 1$), the output grows without bound.
    \item \textbf{Nonlinear phase}---different frequencies are delayed by different
          amounts, which distorts the signal shape.
    \item \textbf{Coefficient sensitivity}---on fixed-point microcontrollers,
          rounding the coefficients can shift poles enough to cause instability.
\end{itemize}

\textbf{Classical IIR design} starts from an analog prototype and maps it to
discrete time:
\begin{itemize}
    \item \textbf{Butterworth}---maximally flat passband, monotonic rolloff. The
          ``vanilla'' choice; sounds boring but is the right default 80\% of the time.
    \item \textbf{Chebyshev Type~I}---equiripple in the passband, monotonic in the
          stopband. Sharper transition than Butterworth for the same order, at the
          cost of passband ripple.
    \item \textbf{Chebyshev Type~II}---monotonic in the passband, equiripple in the
          stopband. Less common but useful when passband flatness matters more.
    \item \textbf{Elliptic (Cauer)}---equiripple in \textit{both} bands. Gives the
          sharpest possible transition for a given filter order, but with ripple
          everywhere.
\end{itemize}

The two main analog-to-digital mappings are:
\begin{itemize}
    \item \textbf{Bilinear transform}: $z = \dfrac{1 + (sT_s/2)}{1 - (sT_s/2)}$.
          This is the workhorse method. It maps the entire $j\omega$ axis onto the
          unit circle with no aliasing, but introduces \textit{frequency warping}:
          analog frequencies get compressed near the Nyquist frequency.
          You compensate by \textbf{pre-warping} the critical frequencies before
          designing the analog prototype.
    \item \textbf{Impulse invariance}: $h[n] = T_s \cdot h_a(nT_s)$.
          Preserves the time-domain impulse response, but suffers from
          \textit{aliasing} if the analog filter is not sufficiently band-limited.
          Mainly used for low-pass and band-pass designs.
\end{itemize}

\paragraph{Head-to-head comparison.}

\begin{center}
\begin{tabular}{lll}
\toprule
\textbf{Property} & \textbf{FIR} & \textbf{IIR} \\
\midrule
Stability & Always stable & Must verify (poles inside unit circle) \\
Phase & Linear phase achievable & Nonlinear phase \\
Order for sharp cutoff & High (50--200+) & Low (4--10 typical) \\
Computation per sample & $M+1$ multiplies & $M+N+1$, but $M,N$ are small \\
Latency & Higher (longer filter) & Lower \\
Design from analog prototype & No & Yes \\
Coefficient sensitivity & Low & High \\
\bottomrule
\end{tabular}
\end{center}

\paragraph{Practical guidelines.}
\begin{itemize}
    \item Need \textbf{linear phase} (audio, symmetric pulse shaping, communications)?
          Use \textbf{FIR}.
    \item Need a \textbf{sharp cutoff with minimal computation} (real-time control
          loops, resource-constrained MCUs)? Use \textbf{IIR}.
    \item Need \textbf{guaranteed stability} without careful analysis? Use
          \textbf{FIR}.
    \item For higher-order IIR filters, \textbf{always} implement as cascaded
          second-order sections (biquads) to reduce coefficient sensitivity:
\end{itemize}
\begin{equation}
H(z) = \prod_{i=1}^{L}
       \frac{b_{0i} + b_{1i}\,z^{-1} + b_{2i}\,z^{-2}}
            {1   + a_{1i}\,z^{-1} + a_{2i}\,z^{-2}}
\end{equation}

Each section is a ``biquad''---a second-order IIR filter that is simple to
implement and numerically well-behaved. Cascading $L$ biquads gives you a
$2L$-th order filter while keeping coefficient quantization effects local to each
section.

\textbf{Example: a biquad in C}
\begin{verbatim}
    // One biquad section
    float biquad(float x, float *state,
                 float b0, float b1, float b2,
                 float a1, float a2) {
        float y = b0*x + state[0];
        state[0] = b1*x - a1*y + state[1];
        state[1] = b2*x - a2*y;
        return y;
    }
\end{verbatim}

This is the \textbf{transposed direct-form~II} structure---two state variables per
biquad, numerically stable, and the standard choice for floating-point
implementations.

\vspace{0.5em}
\begin{mdframed}[linewidth=1pt, innertopmargin=10pt, innerbottommargin=10pt]
\textbf{Quick Decision Flowchart}

\vspace{0.3em}
\begin{enumerate}[nosep]
    \item Is the application phase-sensitive (e.g., symmetric waveform matters)?
          $\rightarrow$ \textbf{FIR} with symmetric coefficients.
    \item Is real-time computation tight (MCU running a fast control loop)?
          $\rightarrow$ \textbf{IIR} (Butterworth or Chebyshev), implemented as
          cascaded biquads.
    \item Is the filter order moderate and your platform has enough memory?
          $\rightarrow$ Either works; default to \textbf{FIR} for simplicity and
          robustness.
    \item Do you need the sharpest possible transition band for a given order?
          $\rightarrow$ \textbf{IIR} with an Elliptic prototype.
\end{enumerate}

\textbf{Rule of thumb for robotics:} Most sensor-smoothing filters on embedded
systems are low-order IIR (often just 1st or 2nd order Butterworth). FIR filters
shine in audio processing, communications, and any application where linear phase
is non-negotiable.
\end{mdframed}

\subsection{Z-Transform}

Everything we have discussed so far with the Laplace transform lives in the continuous world. But your microcontroller is digital---it works with samples, not continuous signals. The \textbf{Z-transform} is the discrete-time equivalent of the Laplace transform.

The Z-transform of a discrete sequence $x[n]$ is:
\begin{equation}
X(z) = \sum_{n=0}^{\infty} x[n] z^{-n}
\end{equation}

The key correspondence is:
\begin{center}
\begin{tabular}{ll}
\toprule
\textbf{Continuous ($s$-domain)} & \textbf{Discrete ($z$-domain)} \\
\midrule
$s$ (derivative) & $\frac{z-1}{T_s}$ (forward difference) \\
$\frac{1}{s}$ (integral) & $\frac{T_s}{z-1}$ (accumulator) \\
Stable if $\text{Re}(s) < 0$ & Stable if $|z| < 1$ \\
$e^{-aT_s}$ in $s$-domain & $z^{-1}$ is one-sample delay \\
\bottomrule
\end{tabular}
\end{center}

\textbf{Why do we care?} Because when you implement a filter or controller on a microcontroller, you are working in the $z$-domain, whether you realize it or not. That simple line of code:
\[
\texttt{y = alpha * x + (1 - alpha) * y\_prev;}
\]
is actually a $z$-domain transfer function:
\[
H(z) = \frac{\alpha}{1 - (1-\alpha)z^{-1}}
\]

Understanding the Z-transform helps you:
\begin{itemize}
    \item Convert continuous-time filters/controllers to code (covered in Section 5).
    \item Analyze the stability of your discrete-time system.
    \item Understand why your controller might behave differently at different sampling rates.
\end{itemize}

For now, just remember: \textbf{$z^{-1}$ means ``one time step ago.''} That is the single most important thing about the Z-transform.

\textbf{Common Z-transform pairs:}

\begin{center}
\begin{tabular}{lll}
\toprule
\textbf{Discrete signal} & \textbf{$z$-domain} & \textbf{When you see it} \\
\midrule
$\delta[n]$ (unit impulse) & $1$ & Initial kick \\
$u[n]$ (unit step) & $\frac{z}{z-1}$ & Step command \\
$a^n u[n]$ (geometric decay) & $\frac{z}{z-a}$ & Discrete first-order response \\
$x[n-1]$ (one-step delay) & $z^{-1}X(z)$ & Previous sample in code \\
$x[n] - x[n-1]$ (first difference) & $(1 - z^{-1})X(z)$ & Discrete derivative \\
$\sum_{k=0}^{n} x[k]$ (running sum) & $\frac{1}{1 - z^{-1}}X(z)$ & Discrete integral \\
\bottomrule
\end{tabular}
\end{center}

\textbf{The connection between $s$ and $z$:} The exact mapping is $z = e^{sT_s}$. A continuous pole at $s = -a$ maps to a discrete pole at $z = e^{-aT_s}$. This means:
\begin{itemize}
    \item A stable continuous pole ($\text{Re}(s) < 0$) maps to $|z| < 1$ (inside the unit circle). Stability is preserved.
    \item A faster continuous pole (more negative $s$) maps to a smaller $|z|$ (closer to the origin). Very fast poles map near $z = 0$.
    \item An oscillatory continuous pole ($s = \sigma \pm j\omega$) maps to $z = e^{\sigma T_s} e^{\pm j\omega T_s}$, which spirals inside (stable) or outside (unstable) the unit circle.
\end{itemize}

\textbf{Reading a $z$-domain transfer function:} Any $z$-domain transfer function can be directly translated to code by replacing $z^{-1}$ with ``previous value.'' For example:
\[
H(z) = \frac{b_0 + b_1 z^{-1}}{1 + a_1 z^{-1}}
\]
becomes:
\begin{verbatim}
    y = b0 * x + b1 * x_prev - a1 * y_prev;
    x_prev = x;  y_prev = y;
\end{verbatim}
This is called a \textbf{difference equation}, and it is the final form of any digital filter or controller.
